% from ega0-22-fr.tex
\oldpage[0\textsubscript{IV-1}]{201}

$D(ty)=tD(y)=0$ puisque $\Omega^1_A\otimes_AK$ est annulé par $\mathfrak m$.
On a ainsi défini l'homomorphisme (22.4.5.2) cherché, et il est immédiat de
vérifier que l'homomorphisme composé
\[
  \Theta_{B_{(K)}/K}\xrightarrow{\pi'_{B_{(K)}/K}}
  \Omega^1_A\otimes_AK\longrightarrow\Omega^1_K
\]
est encore l'homomorphisme canonique $\pi_{B_{(K)}/K}$ (relatif à $i_0$ et
$j_0$). On a donc aussi un homomorphisme canonique analogue à (22.4.5.3).

Si maintenant $K$ est un \emph{corps}, $A/pA$ est une $\mathbf F_p$-algèbre,
et l'on a la suite exacte (20.5.14)
\[
  0\longrightarrow V/pA\longrightarrow
  \Omega^1_{A/pA}\otimes_{A/pA}K\longrightarrow\Omega^1_K\longrightarrow0
\]
et d'autre part on a la suite exacte (20.5.12.1)
\[
  pA\longrightarrow\Omega^1_A\otimes_A(A/pA)
  \longrightarrow\Omega^1_{A/pA}\longrightarrow0
\]
qui montre, en tensorisant avec $K$, que les $K$-espaces vectoriels
$\Omega^1_A\otimes_AK$ et
$\Omega^1_{A/pA}\otimes_{A/pA}K$ sont isomorphes. L'analogue de (22.4.6.2) est
donc ici un homomorphisme
\begin{equation}
\label{0.22.4.8.4-fr}
  \Xi_{B_{(K)}/K}\longrightarrow V/pA=\mathfrak m/(\mathfrak m^2+pA).
  \tag{22.4.8.4}
\end{equation}

Ceci dit, \emph{le critère (22.4.7), où l'on remplace $V$ et $V'$ par $V/pA$
et $V'/pB$ respectivement, reste valable en supposant seulement que $K$ soit
de caractéristique $p>0$} (autrement dit, que $p\cdot1\in\mathfrak m$ dans
$A$), \emph{et que les $F_i(T_i)$ soient de la forme plus générale
(22.4.4.1)}~: en effet, $B$ est un $A$-module libre (donc fidèlement plat), et
il suffit de reprendre le raisonnement de (22.4.7), en y remplaçant (22.4.6.2)
par (22.4.8.4) et $V''$ par $V''/pA''$.
\end{remark}

\subsection{Algèbres géométriquement régulières et algèbres formellement lisses}
\label{subsection:0.22.5-fr}

\begin{proposition}[22.5.1]
\label{0.22.5.1-fr}
Soient $A$ un anneau local régulier, $A'$ un anneau local contenant $A$ et qui
est une $A$-algèbre finie, $\mathfrak m$ (resp. $\mathfrak m'$) l'idéal maximal
de $A$ (resp. $A'$), $K=A/\mathfrak m$ (resp. $K'=A'/\mathfrak m'$) son corps
résiduel. Pour que $A'$ soit un anneau régulier, il faut et il suffit que l'on
ait
\begin{equation}
\label{0.22.5.1.1-fr}
  \operatorname{rg}_K(\mathfrak m/\mathfrak m^2)
  =\operatorname{rg}_{K'}(\mathfrak m'/\mathfrak m'^2).
  \tag{22.5.1.1}
\end{equation}
\end{proposition}

\begin{proof}
En effet, le premier membre de (22.5.1.1) est $\dim(A)$ (17.1.1) et comme
$A\subset A'$ et que $A'$ est une $A$-algèbre finie,
$\dim(A')=\dim(A)$ (16.1.5)~; l'égalité (22.5.1.1) est donc nécessaire et
suffisante pour que $A'$ soit régulier (17.1.1).
\end{proof}

\begin{remark}[22.5.2]
\label{0.22.5.2-fr}
\begin{enumerate}
  \item[\rm(i)] Soient $A_1=A/\mathfrak m^2$,
  $A'_1=A'\otimes_AA_1=A'/\mathfrak m^2A'$~; comme
  $\mathfrak m^2A'\subset\mathfrak m'^2$,
  $\mathfrak m'_1=\mathfrak m'/\mathfrak m^2A'$ est l'idéal maximal de $A'_1$
  et $\mathfrak m'_1/\mathfrak m_1'^2$ est un $K'$-espace vectoriel isomorphe
  à $\mathfrak m'/\mathfrak m'^2$~; la condition de régularité de $A'$ ne
  dépend donc que de la structure de la $A_1$-algèbre $A'_1$, ce qui va nous
  permettre ci-dessous d'appliquer les résultats préliminaires de (22.4) sur
  les algèbres finies sur des anneaux artiniens.
  \item[\rm(ii)] Il résulte de la démonstration de (22.5.1) et de (16.2.6.1)
  que l'on a en tout cas
  \begin{equation}
  \label{0.22.5.2.1-fr}
    \operatorname{rg}_K(\mathfrak m/\mathfrak m^2)
    \leqslant\operatorname{rg}_{K'}(\mathfrak m'/\mathfrak m'^2).
    \tag{22.5.2.1}
  \end{equation}

  